\documentclass[border=10pt,tikz]{standalone}
\usepackage[T2A]{fontenc}
\usepackage[utf8]{inputenc}
\usepackage[russian]{babel}
\usepackage{tikz}
\usetikzlibrary{positioning,arrows.meta,calc,fit,backgrounds}

\begin{document}
	\begin{tikzpicture}[
		font=\sffamily\small,
		>={Latex[length=2.4mm,width=1.8mm]},
		line width=0.5pt,
		node distance=10mm and 14mm,
		block/.style={
			rectangle, draw=black, fill=white,
			minimum width=42mm, minimum height=14mm,
			align=center, inner sep=4pt
		},
		ext/.style={
			rectangle, draw=black, fill=white,
			double, double distance=0.8pt,
			minimum width=38mm, minimum height=14mm,
			align=center, inner sep=4pt
		},
		store/.style={
			rectangle, draw=black, fill=white,
			minimum width=40mm, minimum height=14mm,
			align=center, inner sep=4pt
		},
		subsys/.style={
			rectangle, draw=black, dashed,
			inner sep=6mm
		},
		arr/.style={->, line width=0.5pt},
		lbl/.style={font=\sffamily\footnotesize, inner sep=1.5pt, fill=white}
		]
		
		% ===== Заголовок =====
		\node[font=\sffamily\bfseries] at (0,2.4)
		{Рисунок 1 -- Компонентная архитектура инструмента миграции на ПКК};
		
		% ===== Уровень 1: вход (конфигурация и БД) =====
		\node[store] (cfg)   at (-4, 0)  {Configuration\\Manager};
		\node[store] (sigdb) at ( 4, 0)  {Signature DB\\(JSON)};
		
		% ===== Уровень 2: сканер =====
		\node[block] (scan) at (0,-3.5) {Project Scanner};
		
		% ===== Уровень 3: анализаторы =====
		\node[block] (regex) at (-7,-7.5) {Regex\\Analyzer};
		\node[block] (ast)   at ( 0,-7.5) {AST Analyzer};
		\node[block] (cert)  at ( 7,-7.5) {Cert\\Analyzer};
		
		% ===== Уровень 4: оценщик =====
		\node[block] (risk) at (0,-11.5) {Risk Assessor};
		
		% ===== Уровень 5: генератор =====
		\node[block] (report) at (0,-14.5) {Report Generator};
		
		% ===== Внешние сущности =====
		\node[ext] (user)  at (-13, 0)   {Пользователь\\(CLI)};
		\node[ext] (code)  at (-13,-3.5) {Исходный код\\C/C++};
		\node[ext] (certs) at ( 13,-7.5) {X.509\\сертификаты};
		\node[ext] (json)  at ( 13,-13)  {JSON-отчёт};
		\node[ext] (cbom)  at ( 13,-16)  {CBOM\\CycloneDX 1.5};
		
		% ===== Рамка ядра =====
		\begin{pgfonlayer}{background}
			\node[subsys, fit=(cfg)(sigdb)(scan)(regex)(ast)(cert)(risk)(report)] (core) {};
		\end{pgfonlayer}
		\node[font=\sffamily\bfseries\small, anchor=north west, fill=white, inner sep=2pt]
		at ($(core.south west)+(0.2,-0.05)$) {Ядро системы};
		
		% ===== Связи: внешние входы =====
		\draw[arr] (user.east)  -- node[lbl,above]{параметры} (cfg.west);
		\draw[arr] (code.east)  -- node[lbl,above]{пути} (scan.west);
		% X.509 -> Cert: надпись поднята выше (above=4pt), чтобы не перекрывала стрелку
		\draw[arr] (certs.west) -- node[lbl, above=4pt]{сертификаты} (cert.east);
		
		% Конфигурация и сигнатуры -> сканер
		\draw[arr] (cfg.south)   -- (scan.north west);
		\draw[arr] (sigdb.south) -- (scan.north east);
		
		% ===== Сканер -> анализаторы =====
		% Cert Analyzer: обычная стрелка (всегда)
		\draw[arr] ($(scan.south east)+(-0.5,0)$) -- ($(cert.north)+(0,1.5)$) -- (cert.north);
		
		% Regex/AST: одна стрелка выходит из scan, разветвляется в точке split, сливается в merge
		% Точка разветвления — под серединой scan, между scan и анализаторами
		\coordinate (split) at ($(scan.south)+(0,-1.5)$);
		% Точка слияния — над серединой risk, между анализаторами и risk
		\coordinate (merge) at ($(risk.north)+(0,1.5)$);
		
		% Одна стрелка от scan до split (без наконечника)
		\draw[line width=0.5pt] (scan.south) -- (split);
		
		% От split расходятся две ветки (без наконечников на разветвлении)
		\draw[line width=0.5pt] (split) -| (regex.north);
		\draw[line width=0.5pt] (split) -- (ast.north);
		
		% Метка «или» между ветками
		\node[lbl, draw=black, rounded corners=2pt, inner sep=2pt]
		at (-3.5, -5.8) {или};
		
		% От анализаторов сходятся в merge (без наконечников)
		\draw[line width=0.5pt] (regex.south) |- (merge);
		\draw[line width=0.5pt] (ast.south)   -- (merge);
		
		% Одна стрелка от merge до risk
		\draw[arr] (merge) -- node[lbl,right]{Finding[]} (risk.north);
		
		% Cert -> risk: отдельная стрелка как прежде
		\draw[arr] (cert.south) |- node[lbl,pos=0.9,above]{CertInfo[]} (risk.east);
		
		% Оценщик -> генератор
		\draw[arr] (risk.south) -- node[lbl,right]{RiskReport} (report.north);
		
		% Генератор -> артефакты
		\draw[arr] (report.east) -- ++(2.0,0) |- (json.west);
		\draw[arr] (report.east) -- ++(2.0,0) |- (cbom.west);
		
		% ===== Легенда (сдвинута правее: x=11.0) =====
		\node[
		font=\sffamily\footnotesize,
		anchor=north west,
		draw=black,
		inner sep=6pt,
		align=left,
		fill=white
		]
		at (11.0, 2.8) {%
			\begin{tabular}{@{}c@{\hspace{6pt}}l@{}}
				\tikz\draw[line width=0.5pt] (0,0) rectangle (0.7,0.35);
				& компонент системы \\[4pt]
				\tikz\draw[line width=0.5pt, double, double distance=0.5pt] (0,0) rectangle (0.7,0.35);
				& внешняя сущность \\[4pt]
				\tikz\draw[line width=0.5pt, dashed] (0,0) rectangle (0.7,0.35);
				& подсистема \\[4pt]
				\tikz\draw[->, >=Latex, line width=0.5pt] (0,0.175) -- (0.7,0.175);
				& поток данных \\
			\end{tabular}
		};
		
	\end{tikzpicture}
\end{document}